\section{Tape computations and recovery relations}\label{sec:model}

\paragraph{The machine and its task.}
We use the deterministic multitape model of the first paper in the series
\cite{PaperOne}. The machine has finite control, a two-way read-only input
tape, $k\ge0$ two-way work tapes over a fixed finite alphabet, and an
append-only binary output. All heads start at zero and move by at most one
cell per transition. A transition writes at the scanned work cells and
appends at most one bit. Reading beyond the input gives an end marker.
The \emph{full motion trace} records the head directions and whether a bit
was appended at each transition, together with termination. The next control state, written symbols and appended bit belong to the
step value \cite[\S2]{PaperEleven}.

A $d\times d$ matrix with $b\ge1$ bits per entry occupies $L=d^2b$ bits in
row-major order. For $0\le i,j<d$ and $0\le t<b$, the input bit at address
$q=(id+j)b+t$ must be appended at output position
$\tpos(q)=(jd+i)b+t$. Fix a machine $M$ and a time coefficient $c$.
Write $G_{d,b}$ for the inputs on which $M$ halts within $c(L+1)$ steps
with this output. At the canonical width $b(d)=\max(1,\lceil\log_2 d\rceil)$,
statement \textnormal{(T)} asks for a single $M,c$ with
$G_{d,b(d)}=\{0,1\}^L$ for every positive $d$.

\paragraph{A code and its paired distances.}
Let $V$ be a finite nonempty coordinate set and $\mathcal A$ a finite
alphabet of size at least two. A relation $N$ assigns to each $q\in V$
a set $N(q)\subseteq V$ of permitted helpers. A nonempty set
$C\subseteq\mathcal A^V$ is a \emph{storage code} for $N$ if two words of
$C$ agreeing on $N(q)$ agree at $q$, for every $q$. Its rate is
$\log_{|\mathcal A|}|C|/|V|$. No linearity is assumed.

The incidence graph has an input port $L_v$ and a recovery port $R_v$
for each $v$. Join $L_v$ to $R_q$ when $v\in N(q)$. The \emph{paired
distance} of $q$ is the undirected graph distance from $L_q$ to $R_q$,
with value $\infty$ for disconnected ports. For example, take two
coordinates that store the same bit and let each use its partner as its
only helper. The two incidence edges are disjoint. Both paired distances
are infinite, although each coordinate is recovered exactly. This graph
measurement concerns the permitted observations at the two ports. The example has rate $1/2$ \cite[\S3]{PaperFive}.

A binary family $(V_i,N_i,C_i)$ is \emph{admissible at cap $\Delta$} if
every port has degree at most $\Delta$, $|V_i|\to\infty$, and for each
fixed integer radius $r$ every paired distance exceeds $r$ for all
sufficiently large $i$. Let $R(\Delta)$ be the supremum, over these
families, of their limiting upper rates. Our constructions give lower
bounds on the limiting \emph{lower} rate. For large $\Delta$, write
\[
 p=4k^2+1,\qquad
 f_j(\Delta)=\left(\frac{\log_2\log_2\log_2\Delta}
                         {\log_2\log_2\Delta}\right)^{1/j}\quad(j>0).
\]
The packet and native bounds below use $j=p$ and $j=2p$, respectively.

\paragraph{How a tape run gives a recovery relation.}
Freeze a motion trace. In the diagram of step dependencies, a step has
at most $k+1$ earlier predecessors, one for the previous control step
and one for the last visit to the scanned cell on each work tape. Missing predecessors
give initial values. The input head supplies the bit at its current
address, or the end marker. Choose a set $\Gamma$ of steps and reveal
their full values. This collection of observations is called a
\emph{genie}. Start at the append of output bit $\tpos(q)$ and follow
the predecessor edges backwards, requiring every vertex after the root to lie outside $\Gamma$. The root
itself may be revealed.
The input addresses encountered form the \emph{cone} of that append.
Once the trace and the revealed values have been fixed, agreement on
these input bits forces agreement at the append, by induction on time.
On correctly transposed inputs, it therefore forces agreement at $q$.
This is the recovery relation used at physical bit scale. The same
recursion with packet inputs and packet records will be used in
\cref{sec:proof} \cite[\S4]{PaperEleven}.

For $k$ work walks with $T\le C(n+1)$ steps, $C\ge1$, $n\ge1$, a cut
budget $\varepsilon>0$ and records of width $\beta\ge1$, the bounded-stack theorem gives a cut with
$\beta|\Gamma|\le\varepsilon n$. Each permitted path then has at most
$D+1$ vertices, where $D=\operatorname{genieDepth}(k,\beta,C,\varepsilon)$
is independent of $n$ \cite[Corollary~O(a)]{PaperFourteen}. Here $n$
counts input units. A unit is one packet in \cref{sec:proof} and one
physical bit in \cref{sec:native}. In both applications the revealed
record values are charged in bits.
